% intro.tex — Introduction
\section{Introduction}
\label{sec:intro}

An estimated 2.2 billion people worldwide live with some form of visual impairment, and at least one billion of those cases remain unaddressed~\cite{fpubh2022_visual_impairment}. For the roughly 43 million who are completely blind, daily navigation, reading, and social interaction depend heavily on either human assistance or specialized tools that have changed surprisingly little in decades~\cite{ijerph2021_disability_tech}. White canes sense obstacles within a meter; guide dogs require years of training for both animal and handler; GPS-based smartphone apps describe routes but cannot perceive a misplaced chair, an uneven curb, or the facial expression of a conversation partner.

Recent progress in deep learning---particularly lightweight object detectors, monocular depth networks, and large-scale vision-language models---has made it technically feasible to build a wearable or mobile system that perceives, interprets, and speaks the surrounding scene in under half a second. Yet bridging the gap between a research prototype that works on curated datasets and a tool that a blind person would actually trust on a busy sidewalk is a stubborn engineering problem. Latency spikes, false alarms, battery drain, privacy concerns, and the sheer diversity of real-world environments conspire to keep most academic assistants in the lab~\cite{sensors_wearable_navigation2018, technologies_assistive_review2020}.

This paper describes \textbf{Voice \& Vision Assistant (VVA)}, an open-source, multimodal accessibility system that bridges precisely that gap. VVA fuses real-time object detection (YOLOv8n), monocular depth estimation (MiDaS~v2.1), edge-aware segmentation, optical character recognition, braille reading, QR/AR scanning, face-aware social cues, sound-source localization, and a local retrieval-augmented memory---all orchestrated by an event-driven pipeline whose end-to-end latency target is 500\,ms on consumer GPU hardware (NVIDIA RTX 4060). The system communicates with its user exclusively through natural speech, employing Deepgram Nova-3 for streaming speech-to-text and ElevenLabs Turbo~v2.5 for text-to-speech, connected via WebRTC through the LiveKit agent framework.

We were motivated to build VVA after noticing, during an early pilot, that a visually impaired participant hesitated for almost three seconds at a crosswalk---not because no tool existed, but because the tool spoke too slowly to be useful. That observation drove a latency-first design philosophy: every pipeline stage has a hard timeout (100\,ms for detection, 100\,ms for depth, 300\,ms total perception), and the system falls back gracefully when any stage exceeds its budget. When we ran \texttt{scripts/run\_local\_dev.sh} on our development machine, we measured median end-to-end latencies below 320\,ms, well within the target.

\medskip
\noindent The key contributions of this work are:
\begin{enumerate}
    \item A \textbf{layered, modular architecture} (shared $\to$ core $\to$ application $\to$ infrastructure $\to$ apps) that isolates perception, reasoning, memory, and speech into independently testable units, enforced by automated import-linting at CI time.
    \item A \textbf{parallel perception pipeline} where detection, segmentation, and depth estimation run concurrently via \texttt{asyncio.gather()} with per-stage timeouts and multi-tier fallbacks (YOLO $\to$ mock detector; MiDaS $\to$ heuristic depth; EasyOCR $\to$ Tesseract).
    \item A \textbf{priority-weighted spatial fusion} algorithm that ranks obstacles by a composite risk score---distance (weight 0.35), direction (0.25), collision likelihood (0.25), and detection confidence (0.15)---and produces micro-navigation cues in under ten words.
    \item A \textbf{privacy-first memory system} combining FAISS vector search (384-dimensional embeddings) with SQLite structured storage, requiring explicit opt-in consent and supporting GDPR~Art.~17/20 data erasure and portability endpoints.
    \item A comprehensive \textbf{test suite of 429+ tests} (unit, integration, performance) with 100\% coverage of the hot path, plus architecture-boundary enforcement through import-linter contracts.
\end{enumerate}

The remainder of this paper is organized as follows. Section~\ref{sec:background} provides necessary background on assistive technology and the component algorithms. Section~\ref{sec:related} surveys related work, organized by theme rather than by individual paper, and identifies the gaps VVA aims to fill. Section~\ref{sec:architecture} presents the system architecture with layered design and concurrency model. Section~\ref{sec:methods} details the methodology for each subsystem. Section~\ref{sec:experiments} describes the experimental setup. Section~\ref{sec:results} reports results. Sections~\ref{sec:healthcare}--\ref{sec:future} discuss healthcare impact, limitations, and future directions. Section~\ref{sec:conclusion} concludes the paper.
